\chapter{Concurrency}
\label{chap:concurrency}
See~\cite{cppreference-thread},~\cite{cppreference-memory-order}.
\newline
C++ includes built-in support for threads, atomic operations, mutual exclusion, 
condition variables, and futures. Deeper formal memory-model material lives in 
Chapter~\ref{chap:memorymodel}; this chapter is the library map plus usage notes.
\newline
%%% Paragraph
\section{Threads}
\label{sec:threads}
A \texttt{std::thread} is the C++ handle for an OS thread of execution. Threads in the 
same process \highlight{share the address space}: that is why they communicate cheaply 
through memory, and why unsynchronized shared writes are data races 
(Chapter~\ref{chap:memorymodel}).
\begin{table}[H]
    \begin{tabularx}{\textwidth}{|l|X|}
        \cline{1-2}
        \href{https://en.cppreference.com/w/cpp/thread/thread}{thread}
        & manages a separate thread \\
        \cline{1-2}
        \href{https://en.cppreference.com/w/cpp/thread/jthread}{jthread} 
        & thread with support for auto-joining and cooperative stop (C++20) \\
        \cline{1-2}   
    \end{tabularx}
\end{table}

\toclessnparagraph{Process vs thread (C++)}
\begin{itemize}
    \item A \textbf{process} is an OS instance of a program with its own address space.
    \item A \textbf{thread} is a unit of execution \emph{inside} a process; threads of the 
          same process share memory (hence data races if unsynchronized).
    \item \texttt{std::thread} must be \texttt{join()}ed or \texttt{detach()}ed before 
          destruction, or \texttt{std::terminate} is called --- a common production footgun.
    \item Prefer \texttt{std::jthread} when you can: it joins on destruction and can 
          receive a stop token (C++20).
\end{itemize}
\textbf{Noise / failure modes:} leaking joinable threads at shutdown; oversubscribing 
cores (context-switch jitter) when a thread pool would do; using threads when a 
single-threaded reactor + queues is simpler for I/O.
\par\noindent\textbf{Interview tip:} join-or-detach before destruction; prefer 
\texttt{jthread} when you can explain the stop token.

\begin{lstlisting}[language=C++]
#include <iostream>
#include <thread>
#include <string>

void task1(std::string msg) {
    std::cout << "task1 says: " << msg << '\n';
}

int main() {
    std::thread t1(task1, "Hello");  // starts running
    // ... other work on this thread ...
    t1.join();                       // wait for completion
}
\end{lstlisting}

\subsection{Functions managing the current thread}
\vspace{-8truemm}
These APIs act on the \emph{calling} thread. On a latency path, sleeping and yielding 
are how you \highlight{donate the core} back to the scheduler --- useful for fairness, 
harmful if you needed busy-poll determinism.
\begin{table}[H]
    \begin{tabularx}{\textwidth}{|l|X|}
        \cline{1-2}
        \href{https://en.cppreference.com/w/cpp/thread/yield}{yield} 
        & hints the implementation to reschedule; may be a no-op; not a lock substitute \\
        \cline{1-2}
        \href{https://en.cppreference.com/w/cpp/thread/get_id}{get\char`_id} 
        & returns the id of the current thread \\
        \cline{1-2}
        \href{https://en.cppreference.com/w/cpp/thread/sleep_for}{sleep\char`_for} 
        & blocks the current thread for at least the given duration (wake-up jitter) \\
        \cline{1-2}
        \href{https://en.cppreference.com/w/cpp/thread/sleep_until}{sleep\char`_until} 
        & blocks until a time point \\
        \cline{1-2}
    \end{tabularx}
\end{table}
\textbf{Noise removed by \emph{not} sleeping on the hot path:} wake-up delay and 
unpredictable reschedule time. Prefer condition variables / queues when you must wait; 
prefer busy poll only on isolated cores (Chapter~\ref{chap:lowlatency}).
\par\noindent\textbf{Interview tip:} \texttt{yield} is a hint, not a lock; 
\texttt{sleep\char`_for} is wake-up jitter you chose.
%%% Subsection

\section{Atomicity}
\subsection{Atomic operations}
Fine-grained concurrent updates without a mutex. Each atomic operation is 
indivisible with respect to other atomic operations on the \emph{same} object. 
Atomic objects are 
\href{https://en.cppreference.com/w/cpp/language/memory_model#Threads_and_data_races}{free of data races}
on that object --- that is the race/UB noise they remove compared to a plain 
\texttt{int} shared by two threads.
\newline
Atomics do \emph{not} automatically make multi-word invariants safe; you still choose 
a \texttt{memory\char`_order} (Section~\ref{subsec:memory-order}) so publishes become 
visible in the order you intend.
\par\noindent\textbf{Interview tip:} atomic on one word $\neq$ atomic on a struct of 
several fields --- that still needs a mutex or a carefully designed protocol.

\begin{lstlisting}[language=C++]
std::atomic<int> counter{0};
// concurrent increments are data-race-free:
counter.fetch_add(1, std::memory_order_relaxed);
\end{lstlisting}
%~ \newpage
%%% Paragraph
\subsection{Atomic types}
\vspace{-8truemm}
\texttt{std::atomic<T>} wraps a value so loads/stores/RMWs are atomic. Prefer it for 
flags, counters, and single-word handoffs. \texttt{atomic\char`_ref} overlays atomic 
ops on an existing non-atomic object (lifetime/alignment rules apply) --- useful when 
you cannot change the storage type.
\begin{table}[H]
    \begin{tabularx}{\textwidth}{|l|X|}
        \cline{1-2}
        \multirow{2}{*}{\href{https://en.cppreference.com/w/cpp/atomic/atomic}{atomic}} 
        & atomic class template and specializations for \\
        & bool, integral, and pointer types \\
        \cline{1-2}
        \href{https://en.cppreference.com/w/cpp/atomic/atomic_ref}{atomic\char`_ref}
        & provides atomic operations on non-atomic objects \\
        \cline{1-2}
    \end{tabularx}
\end{table}
Check \texttt{is\char`_lock\char`_free} when you care: some wide types fall back to 
internal locks, which reintroduces contention noise you thought you avoided.
\par\noindent\textbf{Interview tip:} ask whether the atomic is lock-free on the target 
ABI before claiming a lock-free design.
%%% Paragraph

\subsection{Operations on atomic types}
\vspace{-8truemm}
%~ \begin{table}[H]
    %~ \begin{adjustbox}{width=1.1\columnwidth,left}
    %~ \begin{tabular}{|l|l|}
        %~ \cline{1-2}
        %~ \href{https://en.cppreference.com/w/cpp/atomic/atomic_is_lock_free}{atomic\char`_is\char`_lock\char`_free} 
        %~ & checks if the atomic type's operations are lock-free \\
        %~ \cline{1-2}
        %~ \multirow{2}{*}{\parbox{4cm}{\href{https://en.cppreference.com/w/cpp/atomic/atomic_store}{atomic\char`_store} \\ 
                                     %~ \href{https://en.cppreference.com/w/cpp/atomic/atomic_store}{atomic\char`_store\char`_explicit}}}
                         %~ & atomically replaces the value of the atomic object  \\
                         %~ & with a non-atomic argument \\
        %~ \cline{1-2}

        %~ \multirow{2}{*}{\parbox{4cm}{\href{https://en.cppreference.com/w/cpp/atomic/atomic_load}{atomic\char`_load} \\
                                     %~ \href{https://en.cppreference.com/w/cpp/atomic/atomic_load}{atomic\char`_load\char`_explicit}}}
                       %~ & atomically obtains the value stored \\
                       %~ & in an atomic object \\
       %~ \cline{1-2}
       %~ \multirow{2}{*}{\parbox{4.5cm}{\href{https://en.cppreference.com/w/cpp/atomic/atomic_exchange}{atomic\char`_exchange} \\ 
                                      %~ \href{https://en.cppreference.com/w/cpp/atomic/atomic_exchange}{atomic\char`_exchange\char`_explicit}}}
                       %~ &  atomically replaces the value of the atomic object with non-atomic argument\\
                       %~ &  and returns the old value of the atomic\\
        %~ \cline{1-2}
        %~ \multirow{4}{*}{\parbox{7.5cm}{\href{https://en.cppreference.com/w/cpp/atomic/atomic_compare_exchange}{atomic\char`_compare\char`_exchange\char`_weak} \\
                                       %~ \href{https://en.cppreference.com/w/cpp/atomic/atomic_compare_exchange}{atomic\char`_compare\char`_exchange\char`_weak\char`_explicit} \\ 
                                       %~ \href{https://en.cppreference.com/w/cpp/atomic/atomic_compare_exchange}{atomic\char`_compare\char`_exchange\char`_strong} \\ 
                                       %~ \href{https://en.cppreference.com/w/cpp/atomic/atomic_compare_exchange}{atomic\char`_compare\char`_exchange\char`_strong\char`_explicit}}}
                        %~ & \\
                        %~ & atomically compares the value of the atomic object with non-atomic argument \\
                        %~ & and performs atomic exchange if equal or atomic load if not \\
                        %~ & \\
        %~ \cline{1-2}
        %~ \multirow{2}{*}{\parbox{4cm}{\href{https://en.cppreference.com/w/cpp/atomic/atomic_fetch_add}{atomic\char`_fetch\char`_add} \\ 
                                     %~ \href{https://en.cppreference.com/w/cpp/atomic/atomic_fetch_add}{atomic\char`_fetch\char`_add\char`_explicit}}}
                        %~ & adds a non-atomic value to an atomic object\\
                        %~ & and obtains the previous value of the atomic \\
        %~ \cline{1-2}
        %~ \multirow{2}{*}{\parbox{4cm}{\href{https://en.cppreference.com/w/cpp/atomic/atomic_fetch_sub}{atomic\char`_fetch\char`_sub} \\ 
                                     %~ \href{https://en.cppreference.com/w/cpp/atomic/atomic_fetch_sub}{atomic\char`_fetch\char`_sub\char`_explicit}}}
                        %~ & subtracts a non-atomic value from an atomic object\\
                        %~ & and obtains the previous value of the atomic\\
        %~ \cline{1-2}
        %~ \multirow{2}{*}{\parbox{4cm}{\href{https://en.cppreference.com/w/cpp/atomic/atomic_fetch_and}{atomic\char`_fetch\char`_and} \\
                                     %~ \href{https://en.cppreference.com/w/cpp/atomic/atomic_fetch_and}{atomic\char`_fetch\char`_and\char`_explicit}}}
                        %~ & replaces the atomic object with the result of bitwise AND with a non-atomic argument\\
                        %~ & and obtains the previous value of the atomic\\
        %~ \cline{1-2}
        %~ \multirow{2}{*}{\parbox{4cm}{\href{https://en.cppreference.com/w/cpp/atomic/atomic_fetch_or}{atomic\char`_fetch\char`_or} \\
                                     %~ \href{https://en.cppreference.com/w/cpp/atomic/atomic_fetch_or}{atomic\char`_fetch\char`_or\char`_explicit}}}
                        %~ & replaces the atomic object with the result of bitwise OR with a\\
                        %~ & a non-atomic argument and obtains the previous value of the atomic \\
        %~ \cline{1-2}
        %~ \multirow{2}{*}{\parbox{4cm}{\href{https://en.cppreference.com/w/cpp/atomic/atomic_fetch_xor}{atomic\char`_fetch\char`_xor} \\
                                     %~ \href{https://en.cppreference.com/w/cpp/atomic/atomic_fetch_xor}{atomic\char`_fetch\char`_xor\char`_explicit}}}
                        %~ & blocks the thread until notified and the atomic value changes \\
                        %~ & \\
        %~ \cline{1-2}
        %~ \multirow{2}{*}{\parbox{4cm}{\href{https://en.cppreference.com/w/cpp/atomic/atomic_wait}{atomic\char`_wait} \\
                                     %~ \href{https://en.cppreference.com/w/cpp/atomic/atomic_wait}{atomic\char`_wait\char`_explicit}}}
                        %~ & replaces the atomic object with the result of bitwise XOR with a non-atomic argument \\
                        %~ & and obtains the previous value of the atomic \\
        %~ \cline{1-2}
        %~ \href{https://en.cppreference.com/w/cpp/atomic/atomic_notify_one}{atomic\char`_notify\char`_one} 
        %~ & notifies a thread blocked in \texttt{atomic\char`_wait} \\
        %~ \cline{1-2}
        %~ \href{https://en.cppreference.com/w/cpp/atomic/atomic_notify_all}{atomic\char`_notify\char`_all} 
        %~ & notifies all threads blocked in \texttt{atomic\char`_wait} \\
        %~ \cline{1-2}
    %~ \end{tabular}
    %~ \end{adjustbox}
%~ \end{table}
\begin{table}[H]
    %~ \begin{adjustbox}{width=1.1\columnwidth,left}
    \begin{tabularx}{\textwidth}{|l|X|}
        \cline{1-2}
        \href{https://en.cppreference.com/w/cpp/atomic/atomic_is_lock_free}{atomic\char`_is\char`_lock\char`_free} 
        & checks if the atomic type's operations are lock-free \\
        \cline{1-2}
        \multirow{2}{*}{\parbox{4cm}{\href{https://en.cppreference.com/w/cpp/atomic/atomic_store}{atomic\char`_store} \\ 
                                     \href{https://en.cppreference.com/w/cpp/atomic/atomic_store}{atomic\char`_store\char`_explicit}}}
                         & atomically replaces the value of the atomic object  \\
                         & with a non-atomic argument \\
        \cline{1-2}

        \multirow{2}{*}{\parbox{4cm}{\href{https://en.cppreference.com/w/cpp/atomic/atomic_load}{atomic\char`_load} \\
                                     \href{https://en.cppreference.com/w/cpp/atomic/atomic_load}{atomic\char`_load\char`_explicit}}}
                       & atomically obtains the value stored \\
                       & in an atomic object \\
       \cline{1-2}
       \multirow{2}{*}{\parbox{4.5cm}{\href{https://en.cppreference.com/w/cpp/atomic/atomic_exchange}{atomic\char`_exchange} \\ 
                                      \href{https://en.cppreference.com/w/cpp/atomic/atomic_exchange}{atomic\char`_exchange\char`_explicit}}}
                       &  atomically replaces the value of the atomic object with non-atomic argument\\
                       &  and returns the old value of the atomic\\
        \cline{1-2}
        \multirow{4}{*}{\parbox{7.5cm}{\href{https://en.cppreference.com/w/cpp/atomic/atomic_compare_exchange}{atomic\char`_compare\char`_exchange\char`_weak} \\
                                       \href{https://en.cppreference.com/w/cpp/atomic/atomic_compare_exchange}{atomic\char`_compare\char`_exchange\char`_weak\char`_explicit} \\ 
                                       \href{https://en.cppreference.com/w/cpp/atomic/atomic_compare_exchange}{atomic\char`_compare\char`_exchange\char`_strong} \\ 
                                       \href{https://en.cppreference.com/w/cpp/atomic/atomic_compare_exchange}{atomic\char`_compare\char`_exchange\char`_strong\char`_explicit}}}
                        & \\
                        & atomically compares the value of the atomic object with non-atomic argument \\
                        & and performs atomic exchange if equal or atomic load if not \\
                        & \\
        \cline{1-2}
        \multirow{2}{*}{\parbox{4cm}{\href{https://en.cppreference.com/w/cpp/atomic/atomic_fetch_add}{atomic\char`_fetch\char`_add} \\ 
                                     \href{https://en.cppreference.com/w/cpp/atomic/atomic_fetch_add}{atomic\char`_fetch\char`_add\char`_explicit}}}
                        & adds a non-atomic value to an atomic object\\
                        & and obtains the previous value of the atomic \\
        \cline{1-2}
        \end{tabularx}
        \end{table}
        
        \begin{table}[H]
    %~ \begin{adjustbox}{width=1.1\columnwidth,left}
        \begin{tabularx}{\textwidth}{|l|X|}
        \cline{1-2}
        \multirow{2}{*}{\parbox{4cm}{\href{https://en.cppreference.com/w/cpp/atomic/atomic_fetch_sub}{atomic\char`_fetch\char`_sub} \\ 
                                     \href{https://en.cppreference.com/w/cpp/atomic/atomic_fetch_sub}{atomic\char`_fetch\char`_sub\char`_explicit}}}
                        & subtracts a non-atomic value from an atomic object\\
                        & and obtains the previous value of the atomic\\
        \cline{1-2}
        \multirow{2}{*}{\parbox{4cm}{\href{https://en.cppreference.com/w/cpp/atomic/atomic_fetch_and}{atomic\char`_fetch\char`_and} \\
                                     \href{https://en.cppreference.com/w/cpp/atomic/atomic_fetch_and}{atomic\char`_fetch\char`_and\char`_explicit}}}
                        & replaces the atomic object with the result of bitwise AND with a non-atomic argument\\
                        & and obtains the previous value of the atomic\\
        \cline{1-2}
        \multirow{2}{*}{\parbox{4.5cm}{\href{https://en.cppreference.com/w/cpp/atomic/atomic_fetch_or}{atomic\char`_fetch\char`_or} \\
                                     \href{https://en.cppreference.com/w/cpp/atomic/atomic_fetch_or}{atomic\char`_fetch\char`_or\char`_explicit}}}
                        & replaces the atomic object with the result of bitwise OR with a\\
                        & non-atomic argument and obtains the previous value of the atomic \\
        \cline{1-2}
        \multirow{2}{*}{\parbox{4cm}{\href{https://en.cppreference.com/w/cpp/atomic/atomic_fetch_xor}{atomic\char`_fetch\char`_xor} \\
                                     \href{https://en.cppreference.com/w/cpp/atomic/atomic_fetch_xor}{atomic\char`_fetch\char`_xor\char`_explicit}}}
                        & replaces the atomic object with the result of bitwise XOR with a non-atomic argument \\
                        & and obtains the previous value of the atomic \\
        \cline{1-2}
        \multirow{2}{*}{\parbox{4cm}{\href{https://en.cppreference.com/w/cpp/atomic/atomic_wait}{atomic\char`_wait} \\
                                     \href{https://en.cppreference.com/w/cpp/atomic/atomic_wait}{atomic\char`_wait\char`_explicit}}}
                        & blocks the thread until notified and the atomic value changes \\
                        & \\
        \cline{1-2}
        \href{https://en.cppreference.com/w/cpp/atomic/atomic_notify_one}{atomic\char`_notify\char`_one} 
        & notifies a thread blocked in \texttt{atomic\char`_wait} \\
        \cline{1-2}
        \href{https://en.cppreference.com/w/cpp/atomic/atomic_notify_all}{atomic\char`_notify\char`_all} 
        & notifies all threads blocked in \texttt{atomic\char`_wait} \\
        \cline{1-2}
    \end{tabularx}
    %~ \end{adjustbox}
\end{table}

\newpage
\subsection{Flag type and operations}
\vspace{-8truemm}
\begin{table}[H]
    \begin{tabularx}{\textwidth}{|l|X|}
        \cline{1-2}
        \href{https://en.cppreference.com/w/cpp/atomic/atomic_flag}{atomic\char`_flag} 
              & the lock-free boolean atomic type \\
        \cline{1-2}
        \multirow{2}{*}{\parbox{6cm}{\href{https://en.cppreference.com/w/cpp/atomic/atomic_flag_test_and_set}{atomic\char`_flag\char`_test\char`_and\char`_set} \\
                                     \href{https://en.cppreference.com/w/cpp/atomic/atomic_flag_test_and_set}{atomic\char`_flag\char`_test\char`_and\char`_set\char`_explicit}}}
                        & atomically sets the flag to true and returns its previous \\
                        & value\\
        \cline{1-2}
        \multirow{2}{*}{\parbox{4cm}{\href{https://en.cppreference.com/w/cpp/atomic/atomic_flag_clear}{atomic\char`_flag\char`_clear} \\
                                     \href{https://en.cppreference.com/w/cpp/atomic/atomic_flag_clear}{atomic\char`_flag\char`_clear\char`_explicit}}}
                        & atomically sets the value of the flag to false\\
                        &  \\
        \cline{1-2}
        \multirow{2}{*}{\parbox{4cm}{\href{https://en.cppreference.com/w/cpp/atomic/atomic_flag_test}{atomic\char`_flag\char`_test} \\ 
                                        \href{https://en.cppreference.com/w/cpp/atomic/atomic_flag_test}{atomic\char`_flag\char`_test\char`_explicit}}}
                        & atomically returns the value of the flag \\
                        & \\
        \cline{1-2}
        \multirow{2}{*}{\parbox{4cm}{\href{https://en.cppreference.com/w/cpp/atomic/atomic_flag_wait}{atomic\char`_flag\char`_wait} \\ 
                                     \href{https://en.cppreference.com/w/cpp/atomic/atomic_flag_wait}{atomic\char`_flag\char`_wait\char`_explicit}}}
                        & blocks the thread until notified and the flag changes \\
                        & \\
        \cline{1-2}
        \href{https://en.cppreference.com/w/cpp/atomic/atomic_flag_notify_one}{atomic\char`_flag\char`_notify\char`_one} 
        & notifies a thread blocked in atomic\char`_flag\char`_wait \\
        \cline{1-2}
        \href{https://en.cppreference.com/w/cpp/atomic/atomic_flag_notify_all}{atomic\char`_flag\char`_notify\char`_all} 
        & notifies all threads blocked in atomic\char`_flag\char`_wait \\
        \cline{1-2}
    \end{tabularx}
\end{table}

\subsection{Initialization}
\vspace{-8truemm}
\begin{table}[H]
    \begin{tabularx}{\textwidth}{|l|X|}
        \cline{1-2}
        \href{https://en.cppreference.com/w/cpp/atomic/atomic_init}{atomic\char`_init} 
        & non-atomic initialization of a default-constructed atomic \\
        & object \\
        \cline{1-2}
        \href{https://en.cppreference.com/w/cpp/atomic/ATOMIC_VAR_INIT}{ATOMIC\char`_VAR\char`_INIT} 
        & constant initialization of an atomic variable of static storage duration \\
        \cline{1-2}
        \href{https://en.cppreference.com/w/cpp/atomic/ATOMIC_FLAG_INIT}{ATOMIC\char`_FLAG\char`_INIT} 
        & initializes an \texttt{std::atomic\char`_flag} to false\\
        \cline{1-2}
   \end{tabularx}
\end{table}

\newpage
\subsection{Memory synchronization ordering}
\label{subsec:memory-sync-ordering}
\vspace{-8truemm}
Atomicity alone does not decide \emph{when} other threads see surrounding writes. 
\texttt{memory\char`_order} and fences are how you publish a payload (release) and 
consume it (acquire), or take the blunt \texttt{seq\char`_cst} hammer. Formal 
happens-before material: Chapter~\ref{chap:memorymodel}; choice guide: 
Section~\ref{subsec:memory-order}.
\begin{table}[H]
    \begin{tabularx}{\textwidth}{|l|X|}
         \cline{1-2}
         \href{https://en.cppreference.com/w/cpp/atomic/memory_order}{memory\char`_order} 
         & defines memory ordering constraints for the given atomic operation \\
         \cline{1-2}
         \href{https://en.cppreference.com/w/cpp/atomic/kill_dependency}{kill\char`_dependency} 
         & \small{removes the specified object from the \texttt{memory\char`_order\char`_consume} dependency tree} \\
         \cline{1-2}
         \href{https://en.cppreference.com/w/cpp/atomic/atomic_thread_fence}{atomic\char`_thread\char`_fence} 
         & generic memory order-dependent fence synchronization primitive \\
         \cline{1-2}
         \href{https://en.cppreference.com/w/cpp/atomic/atomic_signal_fence}{atomic\char`_signal\char`_fence} 
         & fence between a thread and a signal handler executed in the same thread \\
         \cline{1-2}
    \end{tabularx}
\end{table}
\textbf{Noise removed by correct ordering:} seeing a flag ``ready'' while the payload 
stores are still invisible (or the opposite reordering bugs on weakly ordered CPUs).
\par\noindent\textbf{Interview tip:} release publishes the payload; acquire consumes it --- 
say that before naming fence APIs.

\subsection{Mutual exclusion}
\label{subsec:mutual-exclusion}
See~\cite{cppreference-mutex}.
\newline
A mutex serializes a critical section: only one owner at a time. That removes data 
races on multi-word invariants atomics alone cannot express --- at the cost of 
possible blocking, priority inversion, and cache-line ping-pong under contention.
\newline
\vspace{-8truemm}
\begin{table}[H]
    \begin{tabularx}{\textwidth}{|l|X|}
        \cline{1-2}
        \href{https://en.cppreference.com/w/cpp/thread/mutex}{mutex} 
        & basic exclusive lock; default choice \\
        \cline{1-2}
        \multirow{2}{*}{\href{https://en.cppreference.com/w/cpp/thread/timed_mutex}{timed\char`_mutex} }
        & exclusive lock with try-lock until a timeout \\
        & (detect / bound wait; do not use as a design default) \\ 
        \cline{1-2}
        \multirow{2}{*}{\href{https://en.cppreference.com/w/cpp/thread/recursive_mutex}{recursive\char`_mutex}} 
        & same thread may lock repeatedly; usually a design smell \\
        & (refactor call graphs instead when possible) \\
        \cline{1-2}
        \multirow{2}{*}{\href{https://en.cppreference.com/w/cpp/thread/recursive_timed_mutex}{recursive\char`_timed\char`_mutex}} 
        & recursive + timeout \\
        & \\
        \cline{1-2}
        \href{https://en.cppreference.com/w/cpp/thread/shared_mutex}{shared\char`_mutex} 
        & many readers or one writer; see below \\
        \cline{1-2}
        \multirow{2}{*}{\href{https://en.cppreference.com/w/cpp/thread/shared_timed_mutex}{shared\char`_timed\char`_mutex}} 
        & shared mutex with timeouts \\
        & \\
        \cline{1-2}
    \end{tabularx}
\end{table}
On a latency-critical core, prefer shrinking the critical section or using a lock-free 
SPSC handoff over adding \texttt{recursive\char`_mutex} / heavy reader--writer machinery.
\par\noindent\textbf{Interview tip:} default to \texttt{mutex}; justify 
\texttt{recursive\char`_mutex} or \texttt{shared\char`_mutex} only with a concrete 
contention story.

\subsection{Generic mutex management}
\vspace{-8truemm}
Always lock with RAII. Raw \texttt{lock()}/\texttt{unlock()} pairs leak the lock on 
exceptions and early returns --- that is the deadlock / stuck-mutex class these 
wrappers remove.
\begin{table}[H]
    \begin{tabularx}{\textwidth}{|l|X|}
        \cline{1-2}
        \href{https://en.cppreference.com/w/cpp/thread/lock_guard}{lock\char`_guard} 
        & simplest scope-based exclusive ownership \\
        \cline{1-2}
        \href{https://en.cppreference.com/w/cpp/thread/scoped_lock}{scoped\char`_lock} 
        & RAII for one or many mutexes; avoids lock-order deadlocks when locking several \\
        \cline{1-2}
        \href{https://en.cppreference.com/w/cpp/thread/unique_lock}{unique\char`_lock} 
        & movable exclusive ownership; required by \texttt{condition\char`_variable::wait} \\
        \cline{1-2}
        \href{https://en.cppreference.com/w/cpp/thread/shared_lock}{shared\char`_lock} 
        & movable shared ownership for \texttt{shared\char`_mutex} \\
        \cline{1-2}
        \multirow{3}{*}{\parbox{4cm}{\href{https://en.cppreference.com/w/cpp/thread/lock_tag_t}{defer\char`_lock\char`_t} \\
                                     \href{https://en.cppreference.com/w/cpp/thread/lock_tag_t}{try\char`_to\char`_lock\char`_t} \\
                                     \href{https://en.cppreference.com/w/cpp/thread/lock_tag_t}{adopt\char`_lock\char`_t}}} 
                        &  \\
                        & tag type used to specify locking strategy  \\
                        &  \\
        \cline{1-2}
        \multirow{3}{*}{\parbox{4cm}{\href{https://en.cppreference.com/w/cpp/thread/lock_tag}{defer\char`_lock} \\
                                     \href{https://en.cppreference.com/w/cpp/thread/lock_tag}{try\char`_to\char`_lock} \\
                                     \href{https://en.cppreference.com/w/cpp/thread/lock_tag}{adopt\char`_lock}}} 
                        &  \\
                        & tag constants used to specify locking strategy  \\
                        &  \\
        \cline{1-2}
    \end{tabularx}
\end{table}
\par\noindent\textbf{Interview tip:} \texttt{lock\char`_guard} by default; 
\texttt{unique\char`_lock} when you need \texttt{wait}; \texttt{scoped\char`_lock} 
when locking more than one mutex.
\newpage
\subsection{Generic locking management}
\vspace{-8truemm}
\begin{table}[H]
    \begin{tabularx}{\textwidth}{|l|X|}
        \cline{1-2}
        \href{https://en.cppreference.com/w/cpp/thread/try_lock}{try\char`_lock} 
        & attempts to obtain ownership of mutexes via repeated calls to \texttt{try\char`_lock} \\
        \cline{1-2}
        \href{https://en.cppreference.com/w/cpp/thread/lock}{lock} 
        & locks specified mutexes, blocks if any are unavailable \\
        \cline{1-2}
    \end{tabularx}
\end{table}

\toclessnparagraph{Mutex vs lock vs semaphore (C++)}
\begin{itemize}
    \item A \textbf{mutex} (\texttt{std::mutex}, \ldots) is the synchronization object 
          threads lock/unlock.
    \item A \textbf{lock wrapper} (\texttt{lock\char`_guard}, \texttt{unique\char`_lock}, 
          \texttt{scoped\char`_lock}) is RAII around a mutex --- prefer these over raw 
          \texttt{lock()}/\texttt{unlock()}.
    \item \texttt{std::mutex} is \emph{not} an OS-wide inter-process lock; for that you 
          need OS primitives / shared-memory named mutexes.
    \item A \textbf{counting semaphore} (\texttt{std::counting\char`_semaphore}, C++20) 
          allows up to $N$ concurrent holders; a binary semaphore is similar to a mutex 
          in spirit but with different API/ownership rules.
    \item Reader/writer access: \texttt{shared\char`_mutex} + \texttt{shared\char`_lock} / 
          \texttt{unique\char`_lock}.
\end{itemize}

\begin{lstlisting}[language=C++]
std::mutex m;
int shared = 0;

void inc() {
    std::lock_guard<std::mutex> lock(m);  // unlocks at end of scope
    ++shared;
}
\end{lstlisting}

\toclessnparagraph{Shared mutex (readers / writers)}
\label{npar:shared-mutex}
See~\cite{cppreference-shared-mutex}.
\newline
\texttt{std::shared\char`_mutex} allows \highlight{many concurrent readers} or 
\highlight{one writer}, not both at once:
\begin{itemize}
    \item Readers take \texttt{std::shared\char`_lock} (shared ownership).
    \item Writers take \texttt{std::unique\char`_lock} (exclusive ownership).
\end{itemize}
Use when reads dominate and the critical section is non-trivial. Skip it when 
writers are frequent or the work under the lock is tiny --- writer starvation 
and lock overhead can make a plain \texttt{std::mutex} faster.

\begin{lstlisting}[language=C++]
#include <mutex>
#include <shared_mutex>

std::shared_mutex rw;
int config = 0;   // read often, written rarely

int read_config() {
    std::shared_lock lock(rw);   // many readers OK
    return config;
}

void set_config(int v) {
    std::unique_lock lock(rw);   // exclusive
    config = v;
}
\end{lstlisting}

\par\noindent\textbf{Interview tip:} say ``shared lock for readers, unique lock for 
writers''; then name the trade-off (writer latency / contention vs parallel reads).

\subsection{Call once}
\vspace{-8truemm}
Lazy one-time initialization across threads without hand-rolling a double-checked 
lock. \texttt{call\char`_once} runs the callable at most once; concurrent callers 
block until it finishes (or observe the completed init). Prefer this over 
\texttt{static} locals when you need an explicit flag, or when initialization must 
not run at dynamic init time.
\begin{table}[H]
    \begin{tabularx}{\textwidth}{|l|X|}
        \cline{1-2}
        \href{https://en.cppreference.com/w/cpp/thread/once_flag}{once\char`_flag} 
        & helper object to ensure that \texttt{call\char`_once} invokes the function only once \\
        \cline{1-2}
        \href{https://en.cppreference.com/w/cpp/thread/call_once}{call\char`_once} 
        & invokes a function only once even if called from multiple threads \\
        \cline{1-2}
    \end{tabularx}
\end{table}

\begin{lstlisting}[language=C++]
std::once_flag flag;
void init() { /* runs at most once across all threads */ }

void use() {
    std::call_once(flag, init);
}
\end{lstlisting}
\par\noindent\textbf{Interview tip:} prefer \texttt{call\char`_once} / function-local 
\texttt{static} over a hand-rolled double-checked lock.

\subsection{Condition variables}
\label{subsec:condition-variables}
See~\cite{cppreference-condition-variable}.
\newline
A \texttt{std::condition\char`_variable} lets a thread \highlight{sleep until a 
condition may be true}, instead of spinning. It is always paired with a mutex 
and a \highlight{predicate} (the actual boolean you care about).
\begin{table}[H]
    \begin{tabularx}{\textwidth}{|l|X|}
         \cline{1-2}
         \href{https://en.cppreference.com/w/cpp/thread/condition_variable}{condition\char`_variable} 
         & provides a condition variable associated with a \texttt{std::unique\char`_lock} \\
         \cline{1-2}
         \href{https://en.cppreference.com/w/cpp/thread/condition_variable_any}{condition\char`_variable\char`_any} 
         & provides a condition variable associated with any lock type \\
         \cline{1-2}
         \multirow{2}{*}{\href{https://en.cppreference.com/w/cpp/thread/notify_all_at_thread_exit}{notify\char`_all\char`_at\char`_thread\char`_exit}} 
         & schedules a call to \texttt{notify\char`_all} to be invoked when this thread is \\
         & completely finished \\
         \cline{1-2}
         \href{https://en.cppreference.com/w/cpp/thread/cv_status}{cv\char`_status} 
         & lists the possible results of timed waits on condition variables \\
         \cline{1-2}
    \end{tabularx}
\end{table}

\toclessnparagraph{Why the wait needs the mutex}
\texttt{wait} atomically unlocks the mutex and blocks; on wake it re-locks before 
returning. That way the predicate is always checked while holding the mutex --- 
no lost wakeups between ``see false'' and ``go to sleep.''

\toclessnparagraph{Predicate loop and spurious wakeups}
The OS / library may wake a waiter \highlight{without} a matching 
\texttt{notify} (\highlight{spurious wakeup}). Also, another thread may change 
the shared state again before you run. Therefore:
\begin{itemize}
    \item Always re-check the predicate after wake.
    \item Prefer \texttt{cv.wait(lock, predicate)} (loops until the predicate is true) 
          over bare \texttt{wait(lock)}.
\end{itemize}

\begin{lstlisting}[language=C++]
std::mutex m;
std::condition_variable cv;
bool ready = false;

void consumer() {
    std::unique_lock<std::mutex> lock(m);
    // Equivalent loop form:
    //   while (!ready) cv.wait(lock);
    cv.wait(lock, []{ return ready; });  // wait until predicate is true
    // ... use data while still holding lock if needed ...
}

void producer() {
    {
        std::lock_guard<std::mutex> lock(m);
        ready = true;          // mutate under the mutex
    }                          // unlock before notify is fine (and common)
    cv.notify_one();
}
\end{lstlisting}

\par\noindent\textbf{Interview tip:} ``mutex + predicate + \texttt{wait} with predicate 
(or \texttt{while}); notify after updating state.'' Spurious wakeups $\Rightarrow$ never 
assume the condition is true just because you woke up.

\newpage
\subsection{Futures}
\vspace{-8truemm}
\begin{table}[H]
    \begin{tabularx}{\textwidth}{|l|X|}
        \cline{1-2}
        \href{https://en.cppreference.com/w/cpp/thread/promise}{promise} 
        & stores a value for asynchronous retrieval \\
        \cline{1-2}
        \href{https://en.cppreference.com/w/cpp/thread/packaged_task}{packaged\char`_task} 
        & packages a function to store its return value for asynchronous retrieval \\
        \cline{1-2}
        \href{https://en.cppreference.com/w/cpp/thread/future}{future} 
        & waits for a value that is set asynchronously \\
        \cline{1-2}
        \href{https://en.cppreference.com/w/cpp/thread/shared_future}{shared\char`_future} 
        & waits for a value (possibly referenced by other futures) that is set asynchronously \\
        \cline{1-2}
        \multirow{2}{*}{\href{https://en.cppreference.com/w/cpp/thread/async}{async}} 
        & runs a function asynchronously (potentially in a new thread) and returns \\
        & a \texttt{std::future} that will hold the result  \\
        \cline{1-2}
        \href{https://en.cppreference.com/w/cpp/thread/launch}{launch} 
        & specifies the launch policy for \texttt{std::async} \\
        \cline{1-2}
        \multirow{2}{*}{\href{https://en.cppreference.com/w/cpp/thread/future_status}{future\char`_status}} 
        & specifies the results of timed waits performed on \\
        & \texttt{std::future} and \texttt{std::shared\char`_future} \\
        \cline{1-2}
    \end{tabularx}
\end{table}

\begin{lstlisting}[language=C++]
auto fut = std::async(std::launch::async, []{ return 42; });
int v = fut.get();   // waits if needed; retrieves the value (or rethrows)
\end{lstlisting}

\subsection{Future errors}
\vspace{-8truemm}
Futures fail loudly when misused: get twice, abandon a shared state, break the 
promise/future pairing. These types name that error category --- useful in logs, 
not something to ``catch and ignore'' on a hot path.
\begin{table}[H]
    \begin{tabularx}{\textwidth}{|l|X|}
        \cline{1-2}
        \href{https://en.cppreference.com/w/cpp/thread/future_error}{future\char`_error} 
        & reports an error related to futures or promises \\
        \cline{1-2}
        \href{https://en.cppreference.com/w/cpp/thread/future_category}{future\char`_category} 
        & identifies the future error category \\
        \cline{1-2}
        \href{https://en.cppreference.com/w/cpp/thread/future_errc}{future\char`_errc} 
        & identifies the future error codes \\
        \cline{1-2}
    \end{tabularx}
\end{table}
\newpage

\subsection{Memory Order Semantics}
\label{subsec:memory-order}
See~\cite{cppreference-memory-order},~\cite{sof-memory-order},~\cite{gnugccwiki-atomicsync}.
\newline
For \texttt{std::memory\char`_order} API listings, see Section~\ref{subsec:memory-sync-ordering} above.
\newline
The default for many atomic ops is \texttt{seq\char`_cst}: safest, and often the 
\highlight{expensive} choice on hot paths (extra fences / store buffers on some CPUs). 
Weaker orders are correct only when you understand the publish/consume pattern; 
x86's strong hardware model can hide bugs that appear on ARM.
\newline
Guidance:
\begin{itemize}
    \item \textbf{\texttt{memory\char`_order\char`_relaxed}} --- counters and statistics; atomicity only, no ordering with other memory.
    \item \textbf{\texttt{memory\char`_order\char`_acquire}} --- consumer side of a handoff; no reads/writes may reorder before the load.
    \item \textbf{\texttt{memory\char`_order\char`_release}} --- producer side of a handoff; prior writes must be visible after the store.
    \item \textbf{\texttt{memory\char`_order\char`_acq\char`_rel}} --- RMW that both publishes and consumes (e.g.\ unlock / flag CAS).
    \item \textbf{\texttt{memory\char`_order\char`_seq\char`_cst}} --- safest default when in doubt; single total order across all seq\_cst ops.
\end{itemize}
\textbf{Noise removed by weakening carefully:} unnecessary barriers on every increment. 
\textbf{Noise introduced by weakening wrongly:} readers seeing torn or stale payloads.
\par\noindent\textbf{Interview tip:} name the pair --- \texttt{release} publish, 
\texttt{acquire} consume; use \texttt{acq\char`_rel} on RMWs that do both; default 
\texttt{seq\char`_cst} only when you cannot justify weaker.
\newline
For formal happens-before and sequential consistency, see Section~\ref{subsec:sequential-consistency} in Chapter~\ref{chap:memorymodel}.
Also~\cite{data-dependency-ordering},~\cite{youtube-memory-order-1},~\cite{youtube-memory-order-2}.
